\documentclass[11pt]{article}
\usepackage[T1]{fontenc}
\usepackage{lmodern}
\usepackage{amsmath,amssymb,amsthm}
\usepackage{booktabs}
\usepackage{graphicx}
\usepackage[hidelinks]{hyperref}
\usepackage[margin=1in]{geometry}

\newtheorem{theorem}{Theorem}
\newtheorem{lemma}[theorem]{Lemma}
\newcommand{\mex}{\operatorname{mex}}

\title{Exact Sprague--Grundy Words for the\\
Three-Move Subtraction Games $\{2,5,c\}$}
\author{Xiaojian Zeng\\
\small PBC School of Finance, Tsinghua University, 100190\\
\small \texttt{xiaojianzeng@foxmail.com}}
\date{}

\begin{document}
\maketitle

\begin{abstract}
For the normal-play subtraction game with move set $\{2,5,c\}$, where
$c>5$ and $c\ne7$, Manabe's pure-periodicity criterion leaves the residue
classes $c\equiv0,4\;(\mathrm{mod}\;7)$ as the only cases of his Conjecture~50 not
covered by proved results.  We determine the complete Sprague--Grundy words
in those two classes.  Writing $A=0011021$, for $K\ge1$ we prove
\[
 g_{7K+4}=A^K0011223(021021A^K)^\infty,
 \qquad(q,p)=(7K+7,7K+6),
\]
and for $K\ge2$ we prove
\[
 g_{7K}=A^K32203103(1A^{K-1}02102102)^\infty,
 \qquad(q,p)=(7K+8,7K+2).
\]
Here $q$ is the minimal preperiod and $p$ the least eventual period.
The proof is a parameter-independent finite mex certificate: one seven-row
interior calculation and finitely many seam rows verify the words, while
cyclic gaps between occurrences of the digit $2$ prove period primitivity.
Together with Manabe's results for the other residues, this proves his
Conjecture~50 for the entire shape $(a,b)=(2,5)$.
\end{abstract}

\section{Statement and context}

For $S_c=\{2,5,c\}$ and $n\ge0$, define
\[
 g_c(n)=\mex\{g_c(n-s):s\in S_c,\ s\le n\},
 \qquad g_c(0)=0,
\]
where $\mex$ denotes the least nonnegative integer outside a finite set.
This is the usual Sprague--Grundy recursion for a normal-play impartial
game~\cite{Grundy1939}; one-dimensional take-away games were treated by
Golomb~\cite{Golomb1966}.
Write $g_c=g_c(0)g_c(1)\cdots$ as a one-sided word.  Its minimal preperiod is
\[
 q_c=\min\{q\ge0:\exists p\ge1\text{ such that }
 g_c(n+p)=g_c(n)\text{ for every }n\ge q\},
\]
and $p_c$ is the least such $p$ at $q=q_c$.

Manabe~\cite{Manabe2026} gives an explicit sufficient criterion for pure
periodicity of three-move subtraction games and conjectures its necessity.
For $(a,b)=(2,5)$, his Theorems~28, 32, and~57 give pure periodicity in the
residue classes $1,2,3,5,6\;(\mathrm{mod}\;7)$; the two remaining classes are $0$ and
$4\;(\mathrm{mod}\;7)$.  Our contribution is the following exact resolution of those
classes.

\begin{theorem}\label{thm:main}
Put $A=0011021$.  For every integer $K\ge1$,
\[
 g_{7K+4}=A^K0011223(021021A^K)^\infty,
 \qquad (q_{7K+4},p_{7K+4})=(7K+7,7K+6).
\]
For every integer $K\ge2$,
\[
 g_{7K}=A^K32203103(1A^{K-1}02102102)^\infty,
 \qquad (q_{7K},p_{7K})=(7K+8,7K+2).
\]
Consequently both families have positive minimal preperiod.
\end{theorem}

\section{A finite mex certificate}

Subscripts on a finite word are read cyclically only when explicitly stated.

\begin{lemma}\label{lem:certificate}
Let $X$ be a proposed Sprague--Grundy word with a prefix of length $q$
followed by repetitions of a word $P$ of length $p=c+2$.  To prove that
$X=g_c$, it is enough to check the game recursion in the ranges
\[
 0\le n<c,\qquad c\le n<c+q,
\]
and to check, for every $r\;(\mathrm{mod}\;p)$, the cyclic identity
\begin{equation}\label{eq:cyclic}
 P_r=\mex\{P_{r-5},P_{r-2},P_{r+2}\}.
\end{equation}
\end{lemma}

\begin{proof}
The first two ranges check every index below $c+q$.  If $n\ge c+q$, then
$n-5,n-2,n-c$ all lie in the periodic tail.  Since $p=c+2$, the index
$n-c=n+2-p$ has the same tail phase as $n+2$.  Thus
\eqref{eq:cyclic} is exactly the game recursion for every remaining index.
Forward uniqueness of the mex recursion proves the claim.
\end{proof}

Below $c$, the move $c$ is unavailable.  Direct calculation at
$n=0,\ldots,6$ gives $0011021$.  Thereafter the two-move recursion repeats
$A=0011021$.  The following table records the calculation; the duplicate
third entry is included because $i-5\equiv i+2\;(\mathrm{mod}\;7)$.  In this and all
later tables, each current digit is the mex of the set of digits in its
displayed triple.
\begin{equation}\label{eq:base}
\begin{array}{c|ccccccc}
i&0&1&2&3&4&5&6\\ \hline
A_i&0&0&1&1&0&2&1\\
(A_{i-5},A_{i-2},A_{i+2})&121&111&000&202&111&010&000
\end{array}
\end{equation}

\section{The two missing residue classes}

\subsection{The class c=7K+4}

Set $B=0011223$, $C=021021$, $q=7K+7$, and $p=7K+6$.
The proposed word is $X=A^KB(CA^K)^\infty$, with period block $P=CA^K$.
For $n<c$, the word is the required prefix of $A^\infty$: the first four
letters of $B$ are $0011$, the first four letters of $A$.

The exceptional cyclic rows of \eqref{eq:cyclic} are the following.
All unlisted rows $11\le r\le p-3$ are interior rows of the $A$-run and
reduce to \eqref{eq:base}, with $i\equiv r-6\;(\mathrm{mod}\;7)$; this interval is empty when $K=1$.
\[
\begin{array}{c|rrrrrrrrrrrrr}
r&0&1&2&3&4&5&6&7&8&9&10&p-2&p-1\\ \hline
P_r&0&2&1&0&2&1&0&0&1&1&0&2&1\\
\text{triple}&121&110&002&221&110&000&221&111&000&202&111&010&002
\end{array}
\]

It remains to check $n=c+m$ for $0\le m<q$.  The predecessor triple is
\begin{equation}\label{eq:transient-triple}
 (X_{c+m-5},X_{c+m-2},X_m).
\end{equation}
For $K\ge2$, the first fourteen rows are
\[
\begin{array}{c|rrrrrrrrrrrrrr}
m&0&1&2&3&4&5&6&7&8&9&10&11&12&13\\ \hline
X_{c+m}&2&2&3&0&2&1&0&2&1&0&0&1&1&0\\
\text{triple}&110&010&021&121&130&202&221&310&000&221&111&000&202&111
\end{array}
\]
For $K=1$, this same table is the complete transient check after replacing
the triples at $m=11$ and $m=13$ by $002$ and $113$, respectively.  For
$K\ge2$, the final seven rows, written $m=7K+t$, are
\[
\begin{array}{c|rrrrrrr}
t&0&1&2&3&4&5&6\\ \hline
X_{c+m}&2&1&0&0&1&1&0\\
\text{triple}&010&000&121&111&002&202&113
\end{array}
\]
The remaining range $14\le m\le7K-1$ is empty for $K=2$; otherwise each
row is \eqref{eq:base} with $i\equiv m-2\;(\mathrm{mod}\;7)$.  Indeed, the current letter is
$A_{m-2}$ and the entries of \eqref{eq:transient-triple} have phases
$i-5,i-2,i+2$.  Lemma~\ref{lem:certificate} now proves the asserted word
for every $K\ge1$.

\subsection{The class c=7K}

Here set $D=32203103$, $E=02102102$, $q=7K+8$, and $p=7K+2$.
The proposed word is $X=A^KD(1A^{K-1}E)^\infty$, with
$P=1A^{K-1}E$.  Its first $c$ letters are $A^K$, so the initial check again
follows from \eqref{eq:base}.

The two cyclic seams of $P$ are
\[
\begin{array}{c|rrrrrr}
r&0&1&2&3&4&5\\ \hline
P_r&1&0&0&1&1&0\\
\text{triple}&000&221&111&000&202&111
\end{array}
\]
and, on writing $r=c-8+t$,
\[
\begin{array}{c|rrrrrrrrrr}
t&0&1&2&3&4&5&6&7&8&9\\ \hline
P_r&2&1&0&2&1&0&2&1&0&2\\
\text{triple}&010&002&121&110&002&221&110&002&221&110
\end{array}.
\]
Every row $6\le r\le c-9$ reduces to \eqref{eq:base}, with
$i\equiv r-1\;(\mathrm{mod}\;7)$.  This range is empty for $K=2$, when (6) and (7) are
adjacent and exhaustive.

For the transient, the first fourteen rows of
\eqref{eq:transient-triple} are
\[
\begin{array}{c|rrrrrrrrrrrrrr}
m&0&1&2&3&4&5&6&7&8&9&10&11&12&13\\ \hline
X_{c+m}&3&2&2&0&3&1&0&3&1&0&0&1&1&0\\
\text{triple}&120&110&031&221&120&302&231&210&000&331&111&000&302&111
\end{array}
\]
and, writing $m=c+t$, the final eight rows are
\[
\begin{array}{c|rrrrrrrr}
t&0&1&2&3&4&5&6&7\\ \hline
X_{c+m}&2&1&0&2&1&0&2&1\\
\text{triple}&013&002&122&110&003&221&110&003
\end{array}.
\]
The intervening range $14\le m\le c-1$ is empty for $K=2$ and otherwise
reduces to \eqref{eq:base}, again with $i\equiv m-2\;(\mathrm{mod}\;7)$.  Thus
Lemma~\ref{lem:certificate} proves the second word of
Theorem~\ref{thm:main}.

\section{Minimal period and preperiod}

We now show that the displayed parameters are minimal, rather than merely
valid.  In $P=CA^K$, the occurrences of the digit $2$ are at positions
\[
1,4,11,18,\ldots,7K+4.
\]
Their cyclic gap word is $3\,7^K\,3$.  In $P=1A^{K-1}E$, the cyclic gap
word is, up to rotation, $7^{K-1}3^3$.  Each gap word has exactly one
nonempty cyclic run of $3$'s and one of $7$'s.  If the underlying period
block were a proper power $W^d$, $d\ge2$, then its cyclic gap word between
the $2$'s would repeat the gap word contributed by $W$ and hence would have
at least two runs of each of the two present gap sizes.  This is impossible.
Both period blocks are therefore primitive, including $K=1$ in the first
family and $K=2$ in the second.

It follows that their right-infinite repetitions have least periods
$7K+6$ and $7K+2$, respectively: any period induces a rotation fixing the
finite block, and primitivity forces that rotation to be trivial.

Finally, in both families the letter at $q-1$ is the terminal $3$ of the
transient block.  One period later its phase has letter $1$ for $c=7K+4$
and letter $2$ for $c=7K$.  Suppose periodicity began at some $q'<q$ with
period $d$.  Restriction to the already established tail would force the
least tail period $p$ to divide $d$.  Periodicity at $q-1$ would then equate
the stated unequal letters, a contradiction.  Hence the asserted $q$ is
minimal in both cases, completing the proof of Theorem~\ref{thm:main}.

\section{Completion of the (2,5) shape}

We briefly identify the exact scope inherited from the source.  Specializing
Manabe's notation to $(a,b)=(2,5)$ gives $a+b=7$, $\delta=3$, $\eta=1$,
and $\epsilon=1$.  His Lemma~14, Theorems~28 and~57, and main Theorem~32
prove pure periodicity in precisely the residue classes
\[
 c\equiv1,2,3,5,6\;(\mathrm{mod}\;7).
\]
The two uncovered classes are $0$ and $4$.  Theorem~\ref{thm:main} covers
every admissible member of them: $c=7K$ with $K\ge2$ and $c=7K+4$ with
$K\ge1$, and it proves positive minimal preperiod in both.  The excluded
value $c=7$ is the additive point, while the only smaller boundary $c=6$
belongs to the source-proved residue-$6$ class.  Thus for every $c>5$,
$c\ne7$, the game is purely periodic exactly in the source's admissible
residue classes.  This proves Conjecture~50 of~\cite{Manabe2026} for the
entire shape $(2,5)$, with the exact words above strengthening the required
non-pure-periodicity conclusion in the two complementary classes.

\section*{Computational and AI assistance disclosure}

An exact-integer Python implementation independently checked the mex
recurrence, primitive period blocks, and preperiod-boundary mismatches for
$K=1,\ldots,1000$ in the $7K+4$ family and $K=2,\ldots,1000$ in the $7K$
family.  These finite computations are diagnostic only; the quantified
claims follow from the symbolic range partitions and tables above.  An AI
assistant helped organize the case tables, run the exact checks, and prepare
the manuscript.  The mathematical proof was subjected to a separate audit,
but neither the computation nor the audit is presented as a substitute for
the displayed proof.  This separation between a human-readable reduction
and exact diagnostic verification also appears in finite discrete-geometry
certificates based on integer predicates and independent
implementations~\cite{ZengEtAl2026Certificates}; that work is methodological
context only and is not used in our proof.

\bibliographystyle{plain}
\bibliography{references}

\end{document}
